\section{Transformer Revolution：跨领域迁移的是什么}

在进入信号处理之前，讲者先解释为什么 Transformer 会被不断移植到新领域。真正可迁移的不是“把 token 换成 waveform”这一句口号，而是三类结构能力：任意位置间的内容寻址、逐层更新的共享表示、以及借助 residual/normalization 训练深层堆叠。Audio 会让这些能力与代价同时放大，因为采样率把一秒信号变成上万位置。

\subsection{模型家族像潮汐，不像终局}

CNN、RNN 与 Transformer 在 slide 中被画成动物进化图，强调研究社区会围绕新范式快速集中资源。读图时应把“fashion”理解为研究注意力与工具链聚集，而不是旧机制立刻失效；事实上，后文的 VQ encoder、WaveNet baseline、front-end layer 与 pooling 都继续使用非 attention 组件。

\lecturefigure{slide-03-transformer-revolution-model-zoo.jpg}{模型家族的潮汐：CNN、RNN 与 Transformer 的研究重心迁移。}{Stanford Online 官方视频 00:01:45--00:02:35。}

\teachervoice{讲者说模型族群会“流行”又“退潮”。这个玩笑很重要：本讲会展示 Transformer 获胜的几个表格，但课程结论不是“所有声音都应该 token 化并全局 attention”。}

\subsection{每层 attention 都重新决定什么重要}

早期 encoder--decoder 常把 attention 放在编码完成之后；Transformer 把内容寻址放进每一层，使中间表示可以反复聚合远处信息。讲者的直觉是“cascading self-attention with feature learning”：每层一边寻址，一边把任务相关内容写入新的 representation。这个表述有启发性，但 residual stream 并不会真的删除所有“不重要”内容，它只是让某些方向在后续计算中更有影响。

\lecturefigure{slide-04-transformer-architecture.jpg}{从 terminal attention 到逐层 attention：信息寻址进入每一层。}{Stanford Online 官方视频 00:03:00--00:03:50。}

\lecturefigure{slide-05-early-transformer-history.jpg}{从 BiLSTM-with-attention 到多层自注意力堆叠的历史过渡。}{Stanford Online 官方视频 00:03:50--00:04:28。}

两张图应连续读：第一张回答“attention 放在哪里”，第二张回答“为什么长序列需要更短的信息路径”。RNN 把远距离依赖压进递归状态；self-attention 允许两个时间位置在一层内直接交互，但代价是显式构造大量 pairwise scores。

\lecturefigure{slide-06-usual-transformer-tricks.jpg}{Multi-head attention、skip connection 与 normalization 构成可复用训练配方。}{Stanford Online 官方视频 00:04:30--00:05:12。}

\begin{knowledgebox}{术语消化：Audio Transformer 的对象层级}
\small
\begin{tabularx}{\textwidth}{L{0.20\textwidth}L{0.27\textwidth}X}
\toprule
术语 & 解决的问题 & 在本讲中的具体含义 \\
\midrule
Waveform sample & 连续声压怎样被数字化？ & 采样率 $f_s$ 下的单个离散时刻幅值；raw synthesis 会将其量化为 256 类。 \\
Receptive field & 一个输出能看到多远的过去？ & WaveNet 由 dilation 固定，attention 由数据相关权重选择，但全局选择成本更高。 \\
Residual stream & 多层怎样共享当前表示？ & 每层把 attention/MLP update 加回共同表示，使局部与长程信息可组合。 \\
Spectrogram & 怎样把时域信号变成时频图？ & 对滑动窗口做 Fourier transform 并堆叠 magnitude；时间与频率分辨率互相制约。 \\
Filterbank & 怎样按一组频带分析声音？ & 可以是 mel/constant-Q 等固定滤波器，也可以由前端网络端到端学习。 \\
Codebook token & 连续 embedding 怎样变成离散符号？ & 选择最近 code vector 的索引，供 Transformer 做 next-code language modeling。 \\
Linear probe & 表示是否已含任务信息？ & 冻结 encoder，仅训练线性分类头；它测可线性读取性，不测所有下游能力。 \\
\bottomrule
\end{tabularx}
\end{knowledgebox}

\teachervoice{Verma 特别强调 multi-head、skip connection 和 normalization 不只服务 Transformer；它们是可以迁移到其他网络的训练与表示技巧。名称更新不代表所有机制都从零开始。}

\subsection{Scaling 既放大能力，也放大资源差距}

模型规模曲线在这里承担背景角色：compute 改善和大机构投入让极简单的预测目标被训练到巨大规模，随后形成可迁移 representation。对 Audio 而言，这个逻辑立刻遇到采样率墙——文本一句话可能几十 tokens，十秒 16 kHz 音频却有十六万个 samples。规模不是独立变量，representation 决定同一段内容需要多少 sequence positions。

\lecturefigure{slide-07-model-scaling.jpg}{模型与计算规模增长：简单目标、更多数据和更多资源共同推动表示学习。}{Stanford Online 官方视频 00:05:03--00:06:17。}

\begin{warningbox}{“更大”之前先问 token 是什么}
若把每个 sample 当一个 token，长度随采样率线性增长，attention pair 数随长度平方增长；若先做 spectrogram patch 或 VQ code，长度下降，但细节与归纳偏置改变。参数规模曲线不能脱离表示粒度解读。
\end{warningbox}

\subsection{本章小结}

跨领域迁移的是内容寻址、深层表示更新和成熟训练配方，而不是无条件的全局 attention。Audio 的高采样率会把 pairwise interaction 的代价放大，因此表示选择从第一步就与系统可行性绑定。

